% Bagian: intuitionistic-logic
% Bab: introduction
% Subbagian: syntax

\documentclass[../../../include/open-logic-section]{subfiles}

\begin{document}

\olfileid[id]{int}{int}{syn}

\olsection{Sintaks Logika Intuisionistik}

Sintaks logika intuisionistik sama dengan sintaks logika proposisional.
Dalam logika proposisional klasik, suatu konektif dapat didefinisikan melalui
konektif lain; misalnya, $!A \lif !B$ dapat didefinisikan sebagai $\lnot !A
\lor !B$, atau $!A \lor !B$ sebagai $\lnot(\lnot !A \land \lnot !B)$. Karena
itu, penyajian logika klasik sering memperkenalkan beberapa konektif sebagai
singkatan bagi definisi-definisi tersebut. Hal ini tidak berlaku dalam logika
intuisionistik, dengan dua pengecualian: $\lnot !A$ dapat---dan sering kali
memang---didefinisikan sebagai singkatan bagi $!A \lif \lfalse$. Dengan
demikian, tentu saja, $\lfalse$ sendiri tidak boleh didefinisikan!{} Selain
itu, sebagaimana dalam logika klasik, $!A \liff !B$ dapat didefinisikan sebagai
$(!A \lif !B) \land (!B \lif !A)$.

!!^{formula}s logika proposisional intuisionistik dibentuk dari
\emph{!!{propositional variable}s} dan konstanta proposisional~$\lfalse$
dengan menggunakan \emph{konektif logis}. Kita mempunyai:

\begin{enumerate}
\item !!^a{denumerable} himpunan~$\PVar$ dari !!{propositional variable}s
  $\Obj p_0$, $\Obj p_1$, \dots
\item Konstanta proposisional untuk !!{falsity}~$\lfalse$.
\item Konektif logis: $\land$ (konjungsi), $\lor$ (disjungsi), $\lif$
  (kondisional).
\item Tanda baca: (, ), dan koma.
\end{enumerate}

\begin{defn}[Formula]
\ollabel{defn:formulas} Himpunan~$\Frm[L_0]$ dari \emph{!!{formula}s}
logika proposisional intuisionistik didefinisikan secara induktif sebagai
berikut:
\begin{enumerate}
\item $\lfalse$ merupakan !!{formula} atomik.

\item Setiap !!{propositional variable}~$\Obj p_i$ merupakan !!{formula}
  atomik.

\item Jika $!A$ dan $!B$ merupakan !!{formula}s, maka $(!A \land !B)$
  merupakan !!a{formula}.

\item Jika $!A$ dan $!B$ merupakan !!{formula}s, maka $(!A \lor !B)$
  merupakan !!a{formula}.

\item Jika $!A$ dan $!B$ merupakan !!{formula}s, maka $(!A \lif !B)$
  merupakan !!a{formula}.

\tagitem{limitClause}{Tidak ada hal lain yang merupakan !!a{formula}.}{}
\end{enumerate}
\end{defn}

Selain konektif primitif yang diperkenalkan di atas, kita juga menggunakan
simbol-simbol \emph{terdefinisi} berikut: $\lnot$ (negasi) dan $\liff$
(!!{biconditional}). Formula yang dibentuk menggunakan operator terdefinisi
dipahami sebagai berikut:
\begin{enumerate}
\item $\lnot !A$ menyingkat $!A \lif \lfalse$.

\item $!A \liff !B$ menyingkat $(!A \lif !B) \land (!B \lif !A)$.
\end{enumerate}

Meskipun $\lnot$ secara resmi diperlakukan sebagai singkatan, kadang-kadang
kita akan memberikan aturan dan klausa eksplisit dalam definisi bagi~$\lnot$
seolah-olah simbol itu primitif. Tujuan utamanya ialah agar kita dapat
merumuskan soal latihan.

\end{document}
